% Source: https://artofproblemsolving.com/community/c3h2431035_500th_celebration_post
% Stable identity: aops-topic-2431035-q10
A number is a multiple of $3$ exactly when its digit sum is. The first number has digit sum $7+4+A+5+2+B+1=19+A+B$, so
\[A+B\equiv-19\equiv2\pmod 3.\]
The second has digit sum $3+2+6+A+B+4+C=15+A+B+C$, so $A+B+C\equiv0\pmod3$, hence
\[C\equiv-(A+B)\equiv-2\equiv1\pmod3.\]
So $C\in\{1,4,7\}$, and each of these occurs: for example $A=B=1$ (so $A+B=2$) with $C=1$, $4$ or $7$ gives $7415211$ and $326114C$, which are multiples of $3$. The possible values of $C$ are $\boxed{1,\ 4,\ 7}$.
